% GENERATED by md2tex.py from ../draft_v5_en/front/40_abbreviations.md -- do not edit by hand.
% Change the Markdown and run `python md2tex.py`; edits here are overwritten.

\addchap{List of abbreviations and symbols}

% layout: longtable, natural, \small, 98 of 107 mm
\begingroup
\small\setlength{\tabcolsep}{4pt}
\setlength{\LTleft}{0pt}\setlength{\LTright}{\fill}\setlength{\LTcapwidth}{\textwidth}
\begin{longtable}{@{}ll@{}}
  \toprule
  \tabhead{Abbreviation} & \tabhead{Meaning} \\
  \midrule
  \endfirsthead
  \multicolumn{2}{@{}l}{\emph{(continued)}} \\
  \toprule
  \tabhead{Abbreviation} & \tabhead{Meaning} \\
  \midrule
  \endhead
  \midrule
  \multicolumn{2}{r@{}}{\emph{(continued on next page)}} \\
  \endfoot
  \bottomrule
  \endlastfoot
  ADAS & Advanced Driver-Assistance System \\
  AI & Artificial Intelligence \\
  AMLAS & Assurance of Machine Learning for Autonomous Systems \\
  ASIL & Automotive Safety Integrity Level \\
  CMDP & Constrained Markov Decision Process \\
  CNN & Convolutional Neural Network \\
  CV & Computer Vision \\
  DR & Domain Randomization \\
  DRL & Deep Reinforcement Learning \\
  ESC & Electronic Speed Controller \\
  GSN & Goal Structuring Notation \\
  HARA & Hazard Analysis and Risk Assessment \\
  IMU & Inertial Measurement Unit \\
  KL & Kullback–\allowbreak{}Leibler (divergence between distributions) \\
  MBSE & Model-Based Systems Engineering \\
  MDP & Markov Decision Process \\
  ML & Machine Learning \\
  MPC & Model Predictive Control \\
  ODD & Operational Design Domain \\
  PD & Proportional-Derivative (controller) \\
  PPO & Proximal Policy Optimization \\
  RC & Radio-Controlled \\
  RL & Reinforcement Learning \\
  ROS2 & Robot Operating System 2 \\
  RTX / PhysX & NVIDIA rendering and physics engines (Isaac Sim) \\
  SAC & Soft Actor-Critic \\
  SAE & SAE International (standards body; J3016 levels) \\
  SB3 & Stable-Baselines3 (RL algorithm library) \\
  SBC & Single-Board Computer \\
  SE4AI & Systems Engineering for Artificial Intelligence \\
  SOTIF & Safety of the Intended Functionality (ISO 21448) \\
  SR & Safety Requirement \\
  SRS & Safety Requirements Specification \\
  STPA & System-Theoretic Process Analysis \\
  TTLC & Time-To-Lane-Crossing \\
  V\&V & Verification and Validation \\
\end{longtable}
\endgroup

\textbf{Framework identifiers.} The work uses a single set of identifiers that are never reused. They are summarised here, and their registers are in Appendices A (hazards), B (requirements), E (cage) and F (traceability): \texttt{H-XX} (hazard), \texttt{SR-XXX} (safety requirement), \texttt{SR-CL-A} / \texttt{SR-CL-B} (criticality class of a requirement), \texttt{C-XX} (cage rule), \texttt{SC-*} (scenario), \texttt{M-*} (metric), \texttt{D-NN} (recorded design decision), \texttt{F-X} / \texttt{G-X} (project phase and gate), \texttt{ODD-n.PARAMETER} (declared parameter of the operational domain), \texttt{TBD-Qn} (open question pending closure).

\textbf{Names of trainings and campaigns.} Figures and appendices use short names for the observation tracks, training runs, checkpoints and scenario campaigns: \texttt{F-track}, \texttt{track E}, \texttt{E-main}, \texttt{GE4-V2}, \texttt{margin022}, \texttt{2-D PPO 550k}, \texttt{sim-to-real v2}, and step counts such as \texttt{550k}. Table~\ref{tab:7.1} (\S\ref{sec:7.2.4}) links them to the descriptive names used in the text.

\textbf{Argument identifiers.} These are different from the ones above: they organise the thread of the text, not the chain of evidence. They are \texttt{SO1}\ldots{}\texttt{SO7} (specific objectives, \S\ref{sec:1.4}), \texttt{H1}\ldots{}\texttt{H3} (hypotheses, \S\ref{sec:1.3}; not to be confused with \texttt{H-01}\ldots{}\texttt{H-12}, which are hazards), \texttt{A1}\ldots{}\texttt{A5} (adaptations to the V-Model, \S\ref{sec:3.4}), \texttt{R1}\ldots{}\texttt{R14} (results, Chapter~\ref{ch:12}), \texttt{T1}\ldots{}\texttt{T7} (lines of future work, Chapter~\ref{ch:12}).

\textbf{Main symbols.}

% layout: longtable, natural, \footnotesize, 98 of 107 mm
\begingroup
\footnotesize\setlength{\tabcolsep}{4pt}
\setlength{\LTleft}{0pt}\setlength{\LTright}{\fill}\setlength{\LTcapwidth}{\textwidth}
\begin{longtable}{@{}lll@{}}
  \toprule
  \tabhead{Symbol} & \tabhead{Meaning} & \tabhead{Unit} \\
  \midrule
  \endfirsthead
  \multicolumn{3}{@{}l}{\emph{(continued)}} \\
  \toprule
  \tabhead{Symbol} & \tabhead{Meaning} & \tabhead{Unit} \\
  \midrule
  \endhead
  \midrule
  \multicolumn{3}{r@{}}{\emph{(continued on next page)}} \\
  \endfoot
  \bottomrule
  \endlastfoot
  \texttt{ey} & Lateral deviation from the lane centre & m \\
  \texttt{epsi}, \texttt{θ} & Heading error with respect to the lane tangent & rad / ° \\
  \texttt{κ} & Curvature of the reference trajectory & m\textsuperscript{$-1$} \\
  \texttt{d\_max} & Hard limit on lateral error (rule \cagerule{01}) & m \\
  \texttt{θ\_max} & Limit on heading error (rule \cagerule{02}) & ° \\
  \texttt{t\_min} & Time-to-lane-crossing threshold (rule \cagerule{03}) & s \\
  \texttt{σ\_θ} & Standard deviation of the heading over a sliding window & ° \\
  \texttt{s} & Arc length travelled along the centre line of the circuit & m \\
  \texttt{Δs} & Arc-length increment between two consecutive cycles & m \\
\end{longtable}
\endgroup
